% =============================================================================
% 02-os-operations-protection-notes.tex — Lecture Notes for Lecture 2:
% Operating-System Operations and System Protection
%
% DERIVED ARTIFACT. Source of truth: Slides/CS502/02-os-operations-protection.tex.
% Generated by /lecture-notes; do not hand-edit content independently of the
% Beamer source without re-running /qa-notes afterward (single-source-of-truth.md).
% Structured as a textbook chapter (numbered sections/figures/examples,
% topic-organized rather than a frame-by-frame transcript).
%
% Anchor reading: Silberschatz, Galvin & Gagne, Operating System Concepts,
%                 8th ed. — Ch. 1 §1.5 (dual-mode operation, timer, protection,
%                 pp. 20-23); Stallings, Operating Systems: Internals and Design
%                 Principles, 6th ed. — Ch. 3 §3.5 (p. 140), §3.6 (p. 143);
%                 Bovet & Cesati, Understanding the Linux Kernel, 3rd ed. —
%                 Ch. 4 (pp. 131-188), Ch. 10 (pp. 398-419).
%
% Compile with: cd Notes/CS502 &&
%   TEXINPUTS=../../Preambles;$TEXINPUTS xelatex -interaction=nonstopmode 02-os-operations-protection-notes.tex
%   BIBINPUTS=../..;$BIBINPUTS bibtex 02-os-operations-protection-notes
%   TEXINPUTS=../../Preambles;$TEXINPUTS xelatex -interaction=nonstopmode 02-os-operations-protection-notes.tex  (x2)
% =============================================================================

\documentclass[11pt]{article}
\usepackage[margin=1in]{geometry}
\input{../../Preambles/header}
\coursecode{CS502}

% Chapter-style numbering: this is Lecture 2's chapter.
\renewcommand{\thesection}{2.\arabic{section}}
\renewcommand{\thefigure}{2.\arabic{figure}}
\renewcommand{\thetable}{2.\arabic{table}}
\newtheorem{example}{Example}[section]

\newenvironment{definitionbox}[1]{%
  \par\noindent\textbf{Definition (#1).}\ \itshape
}{\par\vspace{0.3em}}

\title{Operating-System Operations and System Protection\\\large Lecture Notes --- Lecture 2\\
  \normalsize Dual-Mode Operation, Privileged Instructions, and the Timer}
\author{Auqib Hamid Lone\\[0.3em]
  \emph{Department of Computer Science \& Engineering}, \emph{GCET Ganderbal}}
\date{\today}

\begin{document}
\maketitle

\section{The Control Problem}
\noindent\muted{\textbf{Reading.} Silberschatz et al., Ch.~1 \S1.5, pp.~20--23
\cite{SilberschatzOSC_operating_system_concepts}; Stallings, Ch.~3 \S3.5,
p.~140 \cite{StallingsOS_internals_design_principles}.}

A student logs in to a shared laboratory machine and runs a three-token
program: \texttt{while(1);}. The condition is the constant $1$, so it is always
true; the body is empty; the processor therefore executes the test and the
branch forever. The program performs no input, produces no output, issues no
system call, and never exits. On a machine with one core, two programs cannot
both be ``running'' at the same instant: the CPU executes one instruction
stream at a time. The prediction seems unavoidable --- the machine should
freeze. It does not. The instructor, working from another terminal, can still
list the process and kill it. This chapter explains both facts: why a runaway
program cannot monopolize the machine, and why the kernel can always take the
CPU back.

The explanation is not that the operating system watches from outside. The
operating system (OS) is not a remote controller; it is a program that must run
on the same CPU as the code it supervises. Four questions organize what
follows. How does the CPU know who is running? The \key{mode bit} answers it
(Section~2.2). What may user code not do? \key{Privileged instructions} answer
it. How does the OS get the CPU back? The \key{timer} and preemption answer it
(Section~2.3). How does a program ask for help? The \key{system-call gate}
answers it (Section~2.4). Every answer is a hardware mechanism, not a software
convention: protection only works because the hardware enforces it.

\subsection{The OS Is a Program Too}

The kernel uses the same CPU, the same registers, and the same physical memory
as every user program. It cannot run ``in the background'' while user code
runs, because the CPU executes one instruction stream at a time. The kernel
therefore runs only when an \key{event} transfers control to it: a system call,
an interrupt, or an exception. Between events it is not executing at all.

That fact creates the central design problem. The hardware must be able to
tell, at every instant, whether the current instruction stream is
\emph{trusted} kernel code or \emph{untrusted} user code. It must run untrusted
code at full speed, yet guarantee that the trusted kernel can always regain
control:

\begin{quote}
\itshape
Build a machine that can run untrusted code at full speed, yet guarantee that
the trusted kernel can always regain control.
\end{quote}

Every mechanism in this chapter is an answer to that sentence.

\subsection{What Goes Wrong Without a Boundary}

Without a boundary, four failures follow immediately. A buggy pointer writes
over another process's memory --- or over the kernel itself --- because nothing
checks the address. A malicious program disables interrupts, and the timer can
no longer fire. A program reprograms the timer never to expire, then loops
forever. A program issues raw disk commands, bypassing the file system and
every access-control check.

The lesson is blunt: \bad{no boundary means no isolation, no multiplexing, and
no trust.} The rest of this chapter is the mechanism that creates the boundary.

\section{Dual-Mode Operation}
\noindent\muted{\textbf{Reading.} Silberschatz et al., Ch.~1 \S1.5, pp.~20--23
\cite{SilberschatzOSC_operating_system_concepts}.}

Protection begins with a single hardware bit. The bit records which of two
worlds the CPU is currently executing in, and the instruction set itself
changes depending on its value.

\begin{definitionbox}{Dual-mode operation}
A hardware \key{mode bit} records the current privilege level. The kernel runs
in \key{kernel mode} (mode bit $=0$); user programs run in \key{user mode}
(mode bit $=1$). Some instructions are legal only when the bit is $0$.
\end{definitionbox}

The two modes differ in what they may execute and touch. Kernel mode carries
the full instruction set, access to all memory and all devices, and the ability
to change the mode bit. User mode carries a restricted instruction set, access
only to the process's own address space, and no direct way to change the bit.
The restriction is not a convention that well-behaved programs respect; it is
enforced by the hardware on every instruction. The mode bit is \emph{hardware
state}, not an ordinary memory variable \cite{SilberschatzOSC_operating_system_concepts}.

\subsection{Who May Execute What?}

The instruction set is partitioned accordingly. Table~\ref{tab:privileged}
lists the two halves.

\begin{table}[h]
\centering
\small
\begin{tabular}{p{0.44\textwidth}p{0.44\textwidth}}
\toprule
\textbf{Privileged (kernel only)} & \textbf{Unprivileged (user or kernel)} \\
\midrule
\textbullet~set or clear the mode bit & \textbullet~arithmetic and logic \\
\textbullet~disable or enable interrupts & \textbullet~loads and stores to the process's own memory \\
\textbullet~load the timer & \textbullet~branches and function calls \\
\textbullet~issue raw I/O instructions & \textbullet~read the clock (not set it) \\
\textbullet~change the page-table base & \textbullet~execute a trap instruction \\
\textbullet~halt the machine & \\
\bottomrule
\end{tabular}
\caption{The privilege partition. A privileged instruction executed in user
mode traps to the kernel rather than succeeding silently.}
\label{tab:privileged}
\end{table}

Two entries deserve comment. Loading the timer is privileged because user
control of the timer would make preemption optional: the program could simply
set the next expiry beyond its own lifetime. Changing the page-table base is
privileged because it would otherwise let a process redirect address
translation to another process's frames. The privilege boundary is what makes
those operations safe to expose at all.

The last unprivileged entry is the important one: a user program may execute a
trap instruction. That instruction is the only legal door into the kernel, and
Section~2.4 returns to it. A privileged instruction executed in user mode does
not silently fail. It \bad{traps} to the kernel, which normally terminates the
process.

\subsection{The Mode Transition Graph}

Figure~\ref{fig:mode-graph} shows the complete set of legal transitions between
the two modes.

\begin{figure}[h]
  \centering
  \begin{tikzpicture}[every node/.style={font=\small}]
    \node[dag-node, minimum width=3.4cm, minimum height=1.1cm] (user) at (-3.4,0)
      {User mode\\[2pt]\footnotesize mode bit $=1$};
    \node[dag-node, minimum width=3.4cm, minimum height=1.1cm] (kernel) at (3.4,0)
      {Kernel mode\\[2pt]\footnotesize mode bit $=0$};
    \draw[observed-edge] (user.east) --
      node[above, text=primary-blue] {return-from-trap} (kernel.west);
    \draw[observed-edge] (user.north) to[bend left=28]
      node[above, text=primary-blue] {trap / system call} (kernel.north);
    \draw[observed-edge] (user.south) to[bend right=28]
      node[below, text=primary-gold] {interrupt / exception} (kernel.south);
  \end{tikzpicture}
  \caption{The mode transition graph. Entry into the kernel is always a
  controlled hardware event; the return path restores the saved user context
  and the user mode bit. Cf.\ Silberschatz et al., Ch.~1 \S1.5, pp.~20--23
  \cite{SilberschatzOSC_operating_system_concepts}.}
  \label{fig:mode-graph}
\end{figure}

Entry into the kernel is always a \key{controlled} hardware event. A trap or
system call is deliberate; an interrupt or exception is forced. In both cases
the hardware, not user code, performs the mode switch and selects the entry
point. The return path restores the saved user context and the user mode bit.
The kernel does not ``call'' user code; it returns to it.

\subsection{The Mode Bit Is Hardware, Not a Variable}

The bit must be hardware for a simple reason. If it were stored in ordinary
memory, a user program could simply write \texttt{0} into it and become the
kernel. The instruction that changes the bit is itself privileged, so user code
cannot execute it. The hardware defines exactly which events may switch modes:
interrupts, exceptions, and traps. This small, hardware-enforced boundary is
the foundation of the \key{trusted computing base} --- the set of components
the system must trust.

The question is worth asking directly, because it exposes the whole design: why
can't the mode bit live in ordinary memory, and what attack becomes possible if
it does? Suppose the bit were a word in the user's address space. The program
stores \texttt{0} into it. It can then disable interrupts, reprogram the timer,
or overwrite the page tables. The hardware would no longer have a trusted state
to return to. The bit must therefore live outside the address space that user
code can write.

\begin{quote}
\textbf{Consequence.} Protection is not a policy the OS can choose to enforce;
it is a property the hardware guarantees. The OS decides \emph{what} to
protect, but the mode bit decides \emph{whether} it can.
\end{quote}

\subsection{Worked Trace: One \texttt{write} Crosses the Boundary}

Example~\ref{ex:write} follows a single library call across the boundary and
back.

\begin{example}[One \texttt{write} Crosses the Boundary]
\label{ex:write}
User code calls the C library \texttt{write(1, buf, n)}.
\begin{enumerate}
  \item The library places the system-call number and the arguments in
    registers.
  \item It executes a \key{trap} instruction (\texttt{syscall} on x86-64).
  \item Hardware switches to kernel mode, saves the user program counter (PC)
    and flags, and jumps to the trap handler.
  \item The kernel validates the arguments, performs the I/O, and stores a
    return value.
  \item A \key{return-from-trap} restores the user context and the user mode
    bit.
  \item The library returns the result; on failure it sets \texttt{errno}, the
    C library's error variable.
\end{enumerate}
\end{example}

Two details are easy to miss. First, the mode switch is a hardware event, not a
function call: user code does not choose the kernel's entry address, and it
cannot return to an arbitrary instruction inside the kernel. Second, the return
value and \texttt{errno} are set by the library after the kernel returns; the
kernel reports success or failure through the trap's return value. Lecture~3
opens this path up step by step.

Dual-mode operation protects memory and devices. It does not protect time: a
user program can loop forever without executing a single privileged
instruction. The next section supplies the missing mechanism.

\section{The Timer and Preemption}
\noindent\muted{\textbf{Reading.} Silberschatz et al., Ch.~1 \S1.5, pp.~20--23
\cite{SilberschatzOSC_operating_system_concepts}; Stallings, Ch.~3 \S3.5,
p.~140 \cite{StallingsOS_internals_design_principles}.}

\subsection{The Loop That Never Leaves User Mode}

The runaway program is a different kind of problem from the ones dual-mode
operation solves. \texttt{while(1);} executes only unprivileged instructions.
It never traps, never calls the OS, and never touches a device. Every
instruction it executes is legal in user mode. If nothing external interrupts
it, the kernel never runs again, and the machine is not corrupted --- it is
simply \bad{monopolized}.

The OS needs an event that fires \emph{without} the user program's cooperation.
That event is the \key{timer interrupt}.

\subsection{The Timer Is the External Event}

\begin{definitionbox}{Timer interrupt}
An asynchronous event raised by a hardware timer when its count reaches zero.
It switches the CPU into kernel mode and transfers control to the kernel's
timer handler, regardless of what user code is doing.
\end{definitionbox}

The mechanism has five steps. A hardware timer counts down independently of
the instruction stream. When it reaches zero, it raises a \key{timer interrupt}.
The interrupt switches the CPU to kernel mode and jumps to the kernel's timer
handler. The handler saves the running process's context and calls the
scheduler. Loading the timer is a \key{privileged} instruction, so user code
cannot set it.

The timer is a device, not a software alarm. If user code could reload it,
preemption would be optional --- and a malicious program would opt out.

\subsection{A Quantum on the Timeline}

Figure~\ref{fig:quantum} places one scheduling interval, called a quantum, on a
timeline.

\begin{figure}[h]
  \centering
  \begin{tikzpicture}[x=1cm,y=1cm,every node/.style={font=\small}]
    \draw[->, thick, draw=jet] (0,0) -- (12.4,0) node[right] {time};
    \foreach \x/\lab in {0/0,3/10,6/20,9/30,12/40} {
      \draw[jet] (\x,0.08) -- (\x,-0.08);
      \node[below, text=jet] at (\x,0) {\lab\,ms};
    }
    \draw[fill=light-bg, draw=primary-blue, thick] (0,0.35) rectangle (3,0.85);
    \node at (1.5,0.6) {\texttt{P1}};
    \draw[fill=light-bg, draw=primary-gold, thick] (3,0.35) rectangle (6,0.85);
    \node at (4.5,0.6) {\texttt{P2}};
    \draw[fill=light-bg, draw=primary-blue, thick] (6,0.35) rectangle (9,0.85);
    \node at (7.5,0.6) {\texttt{P1}};
    \draw[fill=light-bg, draw=primary-gold, thick] (9,0.35) rectangle (12,0.85);
    \node at (10.5,0.6) {\texttt{P2}};
    \foreach \x in {3,6,9} {
      \draw[primary-gold, thick] (\x,0.9) -- (\x,1.25);
      \node[above, text=primary-gold] at (\x,1.25) {timer};
    }
  \end{tikzpicture}
  \caption{Two processes sharing one CPU in 10\,ms quanta. The timer fires at
  10, 20, and 30\,ms, even when the running process is still runnable. Cf.\
  Silberschatz et al., Ch.~1 \S1.5, pp.~20--23
  \cite{SilberschatzOSC_operating_system_concepts}.}
  \label{fig:quantum}
\end{figure}

\begin{definitionbox}{Quantum}
The fixed interval of CPU time that the scheduler grants one process before the
timer interrupt forces a re-evaluation. In this chapter each process receives a
$10$~ms quantum.
\end{definitionbox}

Each process receives one $10$~ms quantum. The timer fires at the end of every
quantum, even if the process is still runnable. The scheduler decides who runs
next; the timer only guarantees that the decision happens. In
Figure~\ref{fig:quantum}, \texttt{P1} and \texttt{P2} alternate every $10$~ms,
with timer interrupts at $10$, $20$, and $30$~ms.

\subsection{Timer plus Scheduler Equals Preemption}

\begin{definitionbox}{Preemption}
The kernel forcibly takes the CPU away from a running process at a timer
interrupt, saves its context, and may schedule a different process.
\end{definitionbox}

Cooperative scheduling makes yielding voluntary: a process yields when it
chooses to, and one infinite loop stops the system. Preemptive scheduling makes
yielding unnecessary: the timer fires, and the kernel keeps control of the CPU
whether or not the process wanted to give it up. The difference is not
politeness; it is who holds the power to decide.

The choice of quantum is a trade-off. A large quantum wastes response time: an
interactive process waits longer before it next runs. A tiny quantum wastes CPU
time on context switches. The next example makes the second cost quantitative.

\subsection{Worked Example: What Does a Quantum Cost?}

\begin{example}[The Overhead of a Context Switch]
\label{ex:quantum}
A context switch costs $12\ \mu\text{s}$. The quantum is $Q$. What fraction of
CPU time is spent switching?

Let $Q$ be measured in microseconds. The CPU completes $1/Q$ quanta per
microsecond, hence $10^6/Q$ quanta per second, and one context switch occurs at
each quantum boundary. The switching time per second is therefore
\[
  \left(\frac{10^6}{Q}\right)\,(12\ \mu\text{s})
  = \frac{12\times 10^6}{Q}\ \mu\text{s}.
\]
As a fraction of one second ($10^6\ \mu\text{s}$), the overhead is $12/Q$.
Equivalently,
\[
  \text{overhead} = \frac{12\ \mu\text{s}}{Q},
\]
provided $Q$ is expressed in microseconds.

\begin{center}
\renewcommand{\arraystretch}{1.35}
\begin{tabular}{c|c|c}
  \textbf{Quantum $Q$} & \textbf{Switches per second} & \textbf{Overhead} \\
  \hline
  $100\ \mu\text{s}$ & $10{,}000$ & $120{,}000\ \mu\text{s} = 12\%$ \\
  $1\ \text{ms}$     & $1{,}000$  & $12{,}000\ \mu\text{s} = 1.2\%$ \\
  $10\ \text{ms}$    & $100$      & $1{,}200\ \mu\text{s} = 0.12\%$ \\
\end{tabular}
\end{center}

For $Q=100\ \mu\text{s}$, there are $10{,}000$ switches per second; at
$12\ \mu\text{s}$ each, that is $120{,}000\ \mu\text{s}=0.12$~s, or $12\%$ of
the second. For $Q=1$~ms the same calculation gives $1{,}000$ switches and
$12{,}000\ \mu\text{s}=0.012$~s, or $1.2\%$. For $Q=10$~ms it gives $100$
switches and $1{,}200\ \mu\text{s}=0.0012$~s, or $0.12\%$. The formula is the
point: the fixed $12\ \mu\text{s}$ cost is divided by a larger $Q$ as the
quantum grows. This is why the quantum cannot simply be made as small as
possible.
\end{example}

\subsection{Why Not a One-Microsecond Quantum?}

If a small quantum improves response time, why not set the quantum to
$1\ \mu\text{s}$? At $1\ \mu\text{s}$ the switch overhead alone is
$12\ \mu\text{s}/1\ \mu\text{s}=12=1200\%$: the CPU spends twelve times as long
switching as executing. Two further costs compound the problem. Cache and
translation lookaside buffer (TLB) state are destroyed on every switch, so the
next process starts cold. The scheduler itself consumes CPU time to choose the
next process. Real systems therefore choose a quantum that balances
responsiveness against overhead, often in the $1$--$100$~ms range.

\section{Protection in Depth}
\noindent\muted{\textbf{Reading.} Silberschatz et al., Ch.~1 \S1.5, pp.~20--23
\cite{SilberschatzOSC_operating_system_concepts}; Stallings, Ch.~3 \S3.5,
p.~140, and \S3.6, p.~143 \cite{StallingsOS_internals_design_principles}.}

The mode bit and the timer protect the CPU. A complete boundary also protects
memory and the devices.

\subsection{Three Things the Kernel Protects}

Figure~\ref{fig:three-things} summarizes the three resources and the mechanism
that guards each one.

\begin{figure}[h]
  \centering
  \begin{tikzpicture}[every node/.style={font=\small}]
    \node[dag-node, minimum width=3.2cm, minimum height=1.0cm] (cpu) at (-4.2,0)
      {CPU time};
    \node[dag-node, minimum width=3.2cm, minimum height=1.0cm] (mem) at (0,0)
      {Memory};
    \node[dag-node, minimum width=3.2cm, minimum height=1.0cm] (io) at (4.2,0)
      {I/O devices};
    \node[text=primary-gold] at (-4.2,-1.15) {timer};
    \node[text=primary-gold] at (0,-1.15) {MMU};
    \node[text=primary-gold] at (4.2,-1.15) {privileged I/O};
  \end{tikzpicture}
  \caption{The three resources the kernel protects and the hardware mechanism
  behind each. Cf.\ Silberschatz et al., Ch.~1 \S1.5, pp.~20--23
  \cite{SilberschatzOSC_operating_system_concepts}; Stallings, Ch.~3 \S3.5,
  p.~140 \cite{StallingsOS_internals_design_principles}.}
  \label{fig:three-things}
\end{figure}

CPU time is protected by the timer, which forces the kernel to run. Memory is
protected by the memory-management unit (MMU), which checks every user access
against the process's legal address range. I/O is protected because devices are
reachable only through privileged instructions and kernel drivers.

\subsection{Memory Protection}

Three organizations are common, and all three enforce the same rule at
different granularities. With \textbf{contiguous allocation}, base and limit
registers bound the process, and the hardware checks every address against
them. With \textbf{paging}, the page table maps only the process's legal pages,
and a reference to an unmapped page traps to the kernel. With
\textbf{segmentation}, each segment has its own base and limit, and the
hardware checks the offset against the limit.

The pattern is the same in every case: user code never sees a physical address.
Every access is translated and checked by hardware, and a violation becomes a
\key{trap}, not a silent write. Details belong to Unit~III
(Lectures~20--25).

\subsection{I/O Protection}

All raw I/O instructions are privileged. A user program cannot talk to a disk,
a network card, or the timer directly. The user-level \texttt{read} and
\texttt{write} calls are system calls: the kernel validates the buffer, the
device, and the count. It then invokes the device driver, which programs the
controller. The utility \texttt{cat} never touches the disk controller; it asks
the kernel, and the kernel decides whether the request is legal.

This matters because direct device access would let a process bypass file
permissions, corrupt another process's data, or hang the device for everyone.
Privileging the instruction is what makes the kernel's validation unavoidable.

\subsection{Protection versus Security}

Table~\ref{tab:protection-security} separates two terms that are often used
interchangeably.

\begin{table}[h]
\centering
\renewcommand{\arraystretch}{1.4}
\begin{tabular}{p{0.42\textwidth}p{0.42\textwidth}}
\toprule
\textbf{Protection} & \textbf{Security} \\
\midrule
Internal mechanisms among processes &
External threats and attackers \\
memory isolation, privilege levels, access control &
authentication, malware defense, network firewalls \\
Enforced by the OS and hardware &
Enforced by policy, cryptography, and administrative practice \\
\bottomrule
\end{tabular}
\caption{Protection is the mechanism; security is the policy that uses it.
Cf.\ Stallings, Ch.~3 \S3.6, p.~143
\cite{StallingsOS_internals_design_principles}.}
\label{tab:protection-security}
\end{table}

Protection is a \emph{mechanism}; security is a \emph{policy} that uses the
mechanism. Protection concerns internal mechanisms among processes: memory
isolation, privilege levels, and access control. Security concerns external
threats and attackers: authentication, malware defense, and network firewalls.
Protection is enforced by the OS and hardware; security is enforced by policy,
cryptography, and administrative practice. A perfect mode bit does not help if
the administrator grants every user root access
\cite{StallingsOS_internals_design_principles}.

\subsection{The System-Call Gate}

Figure~\ref{fig:syscall-gate} shows the controlled doorway between the two
modes.

\begin{figure}[h]
  \centering
  \begin{tikzpicture}[every node/.style={font=\small}]
    \node[dag-node, minimum width=3.0cm] (user) at (-4.3,0) {User process};
    \node[dag-node, minimum width=2.6cm] (trap) at (-0.8,0) {Trap};
    \node[dag-node, minimum width=3.0cm] (table) at (3.0,0) {System-call table};
    \node[dag-node, minimum width=3.0cm] (svc) at (3.0,-2.0) {Kernel service};
    \draw[observed-edge] (user) -- (trap);
    \draw[observed-edge] (trap) -- (table);
    \draw[observed-edge] (table) -- (svc);
    \draw[observed-edge] (svc) -- ++(-3.6,0) -| (user);
  \end{tikzpicture}
  \caption{The system-call gate: a trap dispatches through the system-call table
  to a kernel service and returns through the same controlled entry point.
  Cf.\ Silberschatz et al., Ch.~2 \S2.3
  \cite{SilberschatzOSC_operating_system_concepts}; Bovet and Cesati, Ch.~10,
  pp.~398--419 \cite{BovetCesati_understanding_linux_kernel}.}
  \label{fig:syscall-gate}
\end{figure}

The trap is the \key{only legal door} from user mode into the kernel. The
kernel validates every argument and every pointer before acting. The return
path restores user mode; the kernel never returns to arbitrary user code.
Figure~\ref{fig:syscall-gate} shows the path: the user process traps, the trap
dispatches through the system-call table to a kernel service, and the service
returns through the same gate. Lecture~3 opens this door and walks through it.

\subsection{The Legal Path for \texttt{cat}}

As a check, suppose all I/O instructions are privileged and ask how
\texttt{cat} reads a file. The legal path has three steps.
\begin{enumerate}
  \item \texttt{cat} calls the library \texttt{read}, which traps into the
    kernel.
  \item The kernel checks the file descriptor, buffer, and count, then asks the
    device driver to perform the transfer.
  \item The driver programs the controller; the kernel returns the data and
    restores user mode.
\end{enumerate}
Nothing in that sequence requires \texttt{cat} to execute a privileged
instruction. The privilege lives in the kernel; the user program only asks.

\section{Putting It Together}
\noindent\muted{\textbf{Reading.} Silberschatz et al., Ch.~1 \S1.5, pp.~20--23
\cite{SilberschatzOSC_operating_system_concepts}; Stallings, Ch.~3 \S3.5,
p.~140 \cite{StallingsOS_internals_design_principles}.}

Protection is not one mechanism. It is a loop: the timer fires, the kernel
runs, the scheduler chooses, and user code resumes --- or traps when it asks
for something it cannot do.

\subsection{The Full Control Loop}

Figure~\ref{fig:control-loop} draws the loop.

\begin{figure}[h]
  \centering
  \begin{tikzpicture}[every node/.style={font=\small}]
    \node[dag-node, minimum width=3.0cm] (user) at (-4.0,-1.0) {User process running};
    \node[dag-node, minimum width=3.0cm] (kernel) at (-4.0,2.0) {Kernel: save context};
    \node[dag-node, minimum width=3.0cm] (sched) at (3.6,2.0) {Scheduler};
    \node[dag-node, minimum width=3.0cm] (restore) at (3.6,-1.0) {Kernel: restore context};
    \draw[observed-edge] (user.north) -- node[left, text=primary-gold] {timer} (kernel.south);
    \draw[observed-edge] (kernel.east) -- (sched.west);
    \draw[observed-edge] (sched.south) -- (restore.north);
    \draw[observed-edge] (restore.west) -- node[below, text=primary-gold] {return-from-trap} (user.east);
  \end{tikzpicture}
  \caption{The full control loop: a timer interrupt transfers control to the
  kernel, the scheduler chooses, and a return-from-trap resumes user code.
  Cf.\ Silberschatz et al., Ch.~1 \S1.5, pp.~20--23
  \cite{SilberschatzOSC_operating_system_concepts}; Stallings, Ch.~3 \S3.5,
  p.~140 \cite{StallingsOS_internals_design_principles}.}
  \label{fig:control-loop}
\end{figure}

The loop runs on every timer tick and every blocking system call. The kernel is
\emph{event-driven}: it runs because something happened, not because it is
``always on''. The user process cannot see the loop; it only sees time passing.

\subsection{Worked Trace: Stopping the Runaway Process}

\begin{example}[Stopping the Runaway Process]
\label{ex:runaway}
\texttt{P1} enters \texttt{while(1);} in user mode at $t=0$. The timeline
unfolds as follows.
\begin{enumerate}
  \item $t=0$: \texttt{P1} enters \texttt{while(1);} in user mode.
  \item $t=10$~ms: the timer fires; hardware switches to kernel mode.
  \item The kernel saves \texttt{P1}'s PC and registers in its process control
    block (PCB).
  \item The scheduler picks \texttt{P2}; the kernel restores \texttt{P2}'s
    context and returns to user mode.
  \item $t=20$~ms: the timer fires again; the scheduler may pick \texttt{P1} or
    a third process.
  \item If \texttt{P1} must be stopped, the kernel can deliver a signal or
    terminate it --- it already has control.
\end{enumerate}
At no step is \texttt{P1} asked to cooperate. The timer is external, the save
is kernel work, and the scheduler's choice is made in kernel mode.
\end{example}

\subsection{Common Misconceptions}

Table~\ref{tab:misconceptions} pairs the four most common wrong statements with
their corrections.

\begin{table}[h]
\centering
\small
\renewcommand{\arraystretch}{1.3}
\begin{tabular}{p{0.42\textwidth}p{0.42\textwidth}}
\toprule
\textbf{Wrong} & \textbf{Right} \\
\midrule
``Dual-mode is a software convention.'' &
The mode bit and privileged instructions are hardware. \\
``The timer is a software alarm.'' &
The timer is a hardware device loaded by the kernel. \\
``User code can disable interrupts if it is fast enough.'' &
Interrupt control is privileged; attempts trap. \\
``The OS runs continuously in the background.'' &
The kernel is event-driven, entered through traps and interrupts. \\
\bottomrule
\end{tabular}
\caption{Four misconceptions about protection, and the hardware fact that
corrects each one.}
\label{tab:misconceptions}
\end{table}

\subsection{Takeaways}

\begin{itemize}
  \item \key{Dual-mode operation} uses a hardware mode bit to separate trusted
    kernel code from untrusted user code.
  \item \key{Privileged instructions} protect the mode bit, the timer, the page
    tables, and the devices.
  \item \key{The timer} is the external event that lets the kernel preempt a
    runaway process.
  \item \key{The system-call gate} is the controlled entry point; the kernel
    validates every request.
  \item Protection is a \emph{hardware-enforced} boundary, not a promise the OS
    makes.
\end{itemize}

\subsection{Looking Ahead: System Calls}

Today's question was how the OS stays in control --- mode bit, privileged
instructions, timer, and the trap gate. The next lecture asks what happens
\emph{inside} the gate: how a program names a service, how parameters travel,
and what kinds of calls exist. The reading is Silberschatz et al., Ch.~2
\S2.3--2.4 \cite{SilberschatzOSC_operating_system_concepts}, and Bovet and
Cesati, Ch.~10 \cite{BovetCesati_understanding_linux_kernel}, for the concrete
Linux view.

\section{Summary}

The operating system is a program that must run on the same CPU as the code it
supervises, so its authority cannot come from standing outside the machine. It
comes from a hardware boundary. A mode bit records whether the CPU is executing
trusted kernel code or untrusted user code, and privileged instructions ---
changing the mode bit, disabling interrupts, loading the timer, issuing raw
I/O, changing the page-table base, halting the machine --- are legal only in
kernel mode. A privileged instruction executed in user mode traps to the kernel
instead of succeeding. Because the bit lives in hardware, outside the address
space user code can write, a program cannot promote itself by storing a value
into memory.

Dual-mode operation protects memory and devices, but not time. A process that
executes only unprivileged instructions can loop forever without ever trapping,
so the kernel needs an event that fires without the process's cooperation. That
event is the timer interrupt: a hardware timer counts down independently,
switches the CPU to kernel mode when it expires, and hands control to the
kernel's timer handler. The handler saves the running process's context, the
scheduler chooses the next process, and a return-from-trap resumes user code.
Preemption is the result: the kernel forcibly takes the CPU away at a timer
interrupt rather than waiting for the process to yield. The quantum controls
the trade-off; with a $12\ \mu\text{s}$ context switch, the overhead is
$12\ \mu\text{s}/Q$, so shrinking the quantum eventually makes switching cost
more than execution.

A complete boundary also covers memory and I/O. Base and limit registers, page
tables, or segment registers ensure that every user access is translated and
checked, and a violation traps rather than writing silently. Raw I/O is
privileged, so user programs reach devices only through system calls that the
kernel validates. Protection is the mechanism; security is the policy that uses
it. The system-call gate is the single controlled entry point, and the kernel
returns to user mode only through the paired return path. In the running
thread, the runaway \texttt{while(1);} process is stopped not by persuading it
to yield but by the timer interrupt, which gives the kernel control and lets it
save the process, schedule another, and terminate the runaway if necessary ---
all without the process's cooperation.

\section*{Self-Check}
\addcontentsline{toc}{section}{Self-Check}

\begin{enumerate}
  \item What does each value of the mode bit mean, and name two instructions
    that are legal only in kernel mode.
  \item Why must the mode bit be hardware rather than an ordinary memory
    variable? Describe the attack that becomes possible if it is not.
  \item A process executes only unprivileged instructions forever. Name the
    mechanism that returns control to the kernel, and outline the sequence.
  \item A context switch costs $12\ \mu\text{s}$. Compute the switching
    overhead for $Q=2$~ms and for $Q=250\ \mu\text{s}$.
  \item Distinguish protection from security, and trace \texttt{cat}'s read in
    three steps.
\end{enumerate}

\subsection*{Exercises}

\begin{enumerate}
  \item A user program executes an instruction that attempts to disable
    interrupts. Describe exactly what happens, using the terms \emph{privileged
    instruction}, \emph{trap}, and \emph{mode bit}. Why is the outcome not the
    same as an ordinary function call?
  \item Classify each event as an interrupt, an exception, or a trap, and state
    which direction it moves the mode bit: (a) the timer expires; (b) a program
    executes \texttt{syscall}; (c) a user program executes a privileged
    instruction. Which of these can the user program prevent?
  \item A context switch costs $12\ \mu\text{s}$. Compute the fraction of CPU
    time spent switching for $Q=250\ \mu\text{s}$ and for $Q=2$~ms. Which
    quantum is the better choice for an interactive system, and what does the
    other cost?
  \item Three processes \texttt{P1}, \texttt{P2}, and \texttt{P3} are ready.
    \texttt{P1} enters \texttt{while(1);} at $t=0$; the scheduler is
    round-robin with a $10$~ms quantum and switch cost is ignored. Describe the
    execution from $t=0$ to $t=40$~ms, and state every instant at which the
    kernel has control and could terminate \texttt{P1}.
  \item A process is confined by base register \texttt{0x10000} and a limit
    register holding the size $64$~KB. Which of these user addresses are legal:
    \texttt{0x10000}, \texttt{0x1FFFF}, \texttt{0x20000}? What happens on an
    illegal access, and why is that outcome preferable to a silent write?
\end{enumerate}

\subsection*{Solutions}
\begingroup\small
\begin{enumerate}
  \item The instruction is privileged, so the hardware checks the current mode
    before executing it. The process is in user mode (mode bit $=1$), so the
    instruction cannot execute. The hardware raises an exception that traps to
    the kernel; the kernel's handler normally terminates the process. The
    outcome differs from a function call because the transfer is forced by
    hardware at a fixed entry point: the user program does not choose the
    kernel address, does not run with kernel privileges, and cannot decide
    whether the kernel returns to it.
  \item (a) The timer expiry is an interrupt: asynchronous and hardware-raised,
    moving the CPU from user mode to kernel mode (or interrupting the kernel
    itself). (b) \texttt{syscall} is a trap: software-initiated by the running
    program, deliberately moving user mode to kernel mode. (c) Executing a
    privileged instruction in user mode is an exception: the hardware rejects
    the instruction and traps to the kernel. The user program can initiate (b)
    but cannot prevent the kernel from handling it; it cannot prevent (a),
    which is external; and it cannot make (c) succeed, because the hardware,
    not the program, decides legality.
  \item With $Q=250\ \mu\text{s}$, there are $1/250\ \mu\text{s}=4{,}000$
    switches per second; at $12\ \mu\text{s}$ each, the switching time is
    $48{,}000\ \mu\text{s}=0.048$~s, or $4.8\%$ of one second. With $Q=2$~ms,
    there are $500$ switches per second, giving
    $6{,}000\ \mu\text{s}=0.006$~s, or $0.6\%$. The $2$~ms quantum is the
    better choice for an interactive system because it keeps overhead at
    $0.6\%$ while still giving each process the CPU frequently; the
    $250\ \mu\text{s}$ quantum improves responsiveness further but spends
    $4.8\%$ of the CPU switching, and also destroys more cache and TLB state.
  \item From $t=0$ to $t=10$~ms, \texttt{P1} runs its infinite loop. At
    $t=10$~ms the timer fires, the kernel saves \texttt{P1}'s context, and the
    scheduler picks \texttt{P2}; \texttt{P2} runs from $10$ to $20$~ms. At
    $t=20$~ms the timer fires, the kernel saves \texttt{P2}, and the scheduler
    picks \texttt{P3}; \texttt{P3} runs from $20$ to $30$~ms. At $t=30$~ms the
    timer fires, and the scheduler picks \texttt{P1} again; \texttt{P1} runs
    from $30$ to $40$~ms, when the next timer interrupt arrives. The kernel
    has control at $t=10$, $20$, $30$, and $40$~ms, so it can deliver a signal
    or terminate \texttt{P1} at any of those instants even though \texttt{P1}
    never yields.
  \item The limit register holds the size, $64$~KB $= \texttt{0x10000}$ bytes,
    so the legal range is $\texttt{0x10000} \le a < \texttt{0x20000}$.
    \texttt{0x10000} is the first legal byte and \texttt{0x1FFFF} is the last;
    both are legal. \texttt{0x20000} is exactly one byte past the end and is
    illegal. The hardware rejects the access and traps to the kernel, which
    normally terminates the process. A silent write would corrupt whichever
    memory happened to lie beyond the limit --- possibly another process's data
    or the kernel itself --- so trapping converts an invisible corruption into
    a controlled, attributable fault.
\end{enumerate}
\endgroup

\subsection*{Further Reading (Page-Cited)}

\begin{itemize}
  \item A.~Silberschatz, P.~B.~Galvin, and G.~Gagne, \emph{Operating System
    Concepts}, 8e --- Ch.~1 \S1.5 (dual-mode operation, timer, protection,
    pp.~20--23).
  \item W.~Stallings, \emph{Operating Systems: Internals and Design
    Principles}, 6e --- Ch.~3 \S3.5 (execution of the OS, p.~140), \S3.6
    (security issues, p.~143).
  \item D.~P.~Bovet and M.~Cesati, \emph{Understanding the Linux Kernel}, 3e
    --- Ch.~4 (interrupts and exceptions, pp.~131--188), Ch.~10 (system calls,
    pp.~398--419), for the concrete x86/Linux view.
\end{itemize}

\noindent\muted{Page-verified via each book's \texttt{index.md} under}\\
\muted{\texttt{master\_supporting\_docs/CS502/}\allowbreak
\texttt{supporting\_books/}.}

\begin{thebibliography}{9}
\bibitem{SilberschatzOSC_operating_system_concepts}
A.~Silberschatz, P.~B.~Galvin, and G.~Gagne, \emph{Operating System Concepts},
8th ed., John Wiley and Sons --- Ch.~1 (pp.~3--47), Ch.~2 (pp.~49--89).
\emph{(Page-cited; index built 2026-10-14.)}
\bibitem{StallingsOS_internals_design_principles}
W.~Stallings, \emph{Operating Systems: Internals and Design Principles}, 6th
ed., Prentice Hall --- Ch.~3 \S3.5 (p.~140), \S3.6 (p.~143).
\emph{(Page-cited; index built 2026-10-14.)}
\bibitem{BovetCesati_understanding_linux_kernel}
D.~P.~Bovet and M.~Cesati, \emph{Understanding the Linux Kernel}, 3rd ed.,
O'Reilly Media, 2006 --- Ch.~4 (interrupts and exceptions, pp.~131--188),
Ch.~10 (system calls, pp.~398--419). \emph{(Page-cited; index built
2026-10-14.)}
\end{thebibliography}

\end{document}
